Tres càrregues elèctriques puntuals, de valor $q = 1,0\ \mathrm{nC}$, es troben situades en els vèrtexs d'un triangle equilàter de 10,0~cm de costat. Dues d'aquestes càrregues són positives, mentre que la tercera és negativa.

\figura[0.3]{fig1.png}

\begin{apartat}{1}
Calculeu la intensitat del camp elèctric en el baricentre del triangle (punt G).
\end{apartat}

\begin{apartat}{1}
Calculeu la variació d'energia potencial electroestàtica que experimenta el sistema si les càrregues se separen fins a formar un altre triangle equilàter de 20,0~cm de costat. Digueu si l'energia potencial electroestàtica augmenta o disminueix i justifiqueu la resposta.
\end{apartat}

\blocetiquetat{Dada:}{$k = \dfrac{1}{4\pi\varepsilon_0} = 8,99\times 10^{9}\ \mathrm{N\,m^2\,C^{-2}}$.}
\nota{El baricentre d'un triangle és el punt d'intersecció de les mitjanes (línies que uneixen cada vèrtex amb el punt mitjà del costat oposat).}
